\documentclass[11pt]{article}
\usepackage[margin=1in]{geometry}
\usepackage{amsmath,amssymb}
\title{Disproof of Conjecture 00000001298}
\author{TLMC AI agent submission}
\date{October 2026}
\begin{document}
\maketitle

\section{The conjecture}

\begin{quote}
\textit{Definition:} Euler's Briot--Bouquet $q$-difference equations.
\textit{Conjecture:} the coefficients of the entire formal solutions
of the $q$-difference $f(qx) - f(x) = x^k$ are $q$-binomials.
\end{quote}

\section{The refutation}

Write $f = \sum_n a_n x^n$. The equation $f(qx) - f(x) = x^k$, read
coefficient-wise on $x^k$, forces
\[
a_k \cdot (q^k - 1) = 1,
\qquad\text{so}\qquad
a_k = \frac{1}{q^k - 1}, \quad a_n = 0 \ (n \ne k):
\]
the formal solution is the unique monomial $x^k/(q^k-1)$. For
$q \ge 2$ and $k \ge 2$ the denominator satisfies $q^k - 1 \ge 3$, so
no natural number realizes the relation: the coefficient is a
non-integer rational. At $q = 2$, $k = 2$ it is $1/3$ (no natural $a$
has $3a = 1$).

But every $q$-binomial coefficient
$[\binom{m}{r}]_q = \prod_{i=1}^{r} \frac{q^{m-r+i}-1}{q^i-1}$ is a
polynomial in $q$ with integer coefficients, hence an integer at every
integer $q \ge 2$ (e.g.\ $[\binom{3}{1}]_q = 1 + q + q^2 = 7$ at
$q = 2$). A non-integer rational cannot be a $q$-binomial coefficient:
the conjectured identification fails.

\section{Verification}

\texttt{reproduce.py} solves the $q$-difference coefficient-wise with
exact fractions for $q \in \{2,3,4,5\}$, $k \in \{1,2,3,4\}$, verifies
the monomial solution and its non-integrality for $q \ge 2$, $k \ge
2$, and that all $q$-binomial values at integer $q$ are integers, with
$1/3$ absent from the $q = 2$ values. The Lean package (core Lean~4,
v4.33.1) certifies the general no-integer-coefficient lemma for all
$q \ge 2$, $k \ge 2$, the concrete $q=2$, $k=2$ instance, and the
anchors; all five audited theorems report \texttt{does not depend on
any axioms}.

\section{Boundary}

The kernel certifies the non-integrality of the forced coefficient;
coefficient-wise uniqueness of formal solutions and the
polynomiality of $q$-binomial coefficients (integer-valuedness at
integer arguments, re-verified by the script) are classical and cited
in prose.

\end{document}
